\subsection{The finite path barrier and the state bound}
\label{sec:path-details}
The reachability proof uses the following three direct identities of
the original rightmost operation. For every legal recursive prefix
$P\spref X$,
\begin{equation}\label{eq:9-3}
 \begin{gathered}
 X[n]\pref X[n+1],\qquad
 P\pref X[n]\ \text{or}\ X[n]\pref P,\qquad
 P\pref X[1].
 \end{gathered}                                                
\end{equation}
To verify them, locate the cut of $P$ on the rightmost active branch.
The general zero case deletes the deepest last column. A positive
case keeps the prefix before its final bare entry and fills only the
rightmost hole; increasing the index fills that same hole further.
An ancestor root does not rewrite the retained prefix. The special
case deletes the singleton owner-bare column and appends copies of
the same shortened enclosing entry; index one keeps at least one
such entry. If index zero removes an enclosing entry entirely, its
output is below the corresponding prefix cut. Empty-last-column
deletion is immediate. Comparing the cut with these positions proves
(\ref{eq:9-3}) by finite recursion on the containing branch.

If $P\pref X$ and a given finite path from $X$ ends at $Y>P$, then
$P\spref Y$: follow (\ref{eq:9-3}) edge by edge. Once a path reaches $P$ or
below, strict native decrease prevents its return above $P$. At the
first crossing below $P$, its output is a recursive prefix of $P$.

\begin{lemma}[Two finite paths]\label{lem:two-paths}
Suppose $X,Y$ are recursive-prefix comparable, with given finite paths
$X\to^*H$ and $Y\to^*B$. If $P=H[n]$ and $P<B<H$, then
$P\spref B$.
\end{lemma}
\begin{proof}
Induct on the sum of the remaining lengths of these two given paths,
keeping $H,B,P$ fixed. If $X=Y=H$, the first child $V$ on the path
to $B$ is above $P$. The children of $H$ are prefix comparable by
(\ref{eq:9-3}), so $P\spref V$, and the preceding finite-path barrier gives
$P\spref B$. If $X=Y\ne H$, neither endpoint is the common starting
point: $B<H\le X$. Take the first edge on both paths; the new starting
points are children of one node and hence prefix comparable. The
remaining total length decreases.

If $X\spref Y$, then $B<H\le X$. On the given path from $Y$ to $B$,
take its first point $Y'\le X$. The barrier gives $Y'\pref X$.
At least one edge has been consumed, so apply the induction hypothesis
to the remaining paths. If $Y\spref X$ and $H>Y$, the path barrier
first gives $Y\spref H$. Since $P<B\le Y$, (\ref{eq:9-3}) at $H\to P$
forces $P\spref Y$, and the path to $B$ gives the conclusion. If
instead $H\le Y$, take the first $X'\le Y$ on the path from $X$
to $H$. Again $X'\pref Y$ and a positive number of edges is consumed.
All unfinished branches strictly reduce the stated natural-number
measure. No well-ordering of the descendant relation is assumed.
\end{proof}

For the size bound, the legal recursive prefixes of a finite $K$ form
a chain of at most $\nu(K)+1$ members: each possible stopping place
corresponds to a column or entry in depth-first order, and invalid empty
private-column stops are omitted. If $K<X$ is not itself a prefix,
let $\rho_K(X)$ be the least legal recursive prefix of $X$ above $K$.
At the first difference of $X$ and $K$, retain just that difference
and truncate everything later. A differing parent needs its entry
with empty height. If $K$ ends first, retain one additional entry,
one additional top-level empty column, or a private column and its
first bare entry. Thus
\begin{equation}\label{eq:9-4}
          K<\rho_K(X)\pref X,\qquad
          \nu(\rho_K(X))\le\nu(K)+2.                       
\end{equation}
The earlier containing framework was already counted in $K$.

Put $N=\nu(K)$. If a compressed state $U$ is a limit type, its
$N+2$ children with indices $0,\ldots,N+1$ are distinct: each
additional general copy adds a nonempty native window, and each
special copy adds a shortened entry. If they were all below $K$
without a rejecting nonprefix crossing, they would be $N+2$ distinct
strict prefixes of $K$, contradicting its prefix bound. For a successor
type, the deletion child is the immediate predecessor, so it is at
least $K$. This proves the finite child bound used in
Lemma~\ref{lem:reachability}.

Finally, every parent in a scope-correct expression of at most $N+2$
objects refers to a visible earlier column; its address is bounded by
the total number of columns. There are finitely many tree shapes and
such address assignments. Compressed states therefore range over a
fixed finite set and strictly decrease, proving termination of the
search. If a rejecting crossing occurred for actual $K$ and $V$, their
given histories start at prefix-comparable towers; apply
Lemma~\ref{lem:two-paths} to that crossing to contradict rejection.
This establishes reachability without using any ordinal labels.
